\section{Architettura del sistema}
\label{sec:architettura}

\subsection{Il ciclo sense-think-act-log}

Il filo conduttore dell'intera architettura è un ciclo continuo, ispirato al
paradigma degli agenti razionali e dei sistemi robotici/IoT, articolato in
quattro fasi che si ripetono indefinitamente:

\begin{itemize}
    \item \textbf{SENSE} (percepisci): un sensore visivo virtuale acquisisce
    un'immagine di una foglia, un sensore ambientale virtuale rileva
    temperatura, umidità dell'aria, umidità del suolo e luminosità.
    \item \textbf{THINK} (pensa): la CNN classifica l'immagine, l'agente
    decisionale combina la diagnosi con i dati ambientali e decide le azioni
    da intraprendere.
    \item \textbf{ACT} (agisci): gli attuatori virtuali (pompa di
    irrigazione, ventola, allarme, notifica) vengono attivati o disattivati
    in base alla decisione.
    \item \textbf{LOG} (registra): l'esito del ciclo (predizione, confidenza,
    dati ambientali, azioni intraprese, metriche di performance) viene
    registrato per l'analisi successiva.
\end{itemize}

La Figura~\ref{fig:diagramma-blocchi} illustra il flusso dei dati fra i
componenti software che realizzano questo ciclo.

\begin{figure}[h]
\centering
\begin{tikzpicture}[
    box/.style={draw, rounded corners, minimum width=3.6cm, minimum height=1.1cm,
        align=center, fill=pomodoroTeal!8, draw=pomodoroTeal, thick, font=\small},
    node distance=1.4cm and 1.8cm,
    arrow/.style={-{Stealth[length=2.5mm]}, thick, pomodoroDark},
]
\node[box] (camera) {Sensor Layer\\\footnotesize Virtual Camera + Env. Simulator};
\node[box, right=of camera] (edge) {Edge AI Layer\\\footnotesize Inference Engine (CNN compressa)};
\node[box, right=of edge] (act) {Actuator Layer\\\footnotesize Irrigazione / Ventola / Allarme};

\node[box, below=of edge, fill=pomodoroPurple!8, draw=pomodoroPurple] (agent) {Agente Decisionale\\\footnotesize Regole PEAS + Data Fusion};
\node[box, below=of agent, fill=pomodoroAmber!10, draw=pomodoroAmber] (log) {Logging / Dashboard};

\draw[arrow] (camera) -- (edge);
\draw[arrow] (edge) -- (act);
\draw[arrow] (edge) -- (agent);
\draw[arrow] (camera) -- (agent);
\draw[arrow] (agent) -- (act);
\draw[arrow] (agent) -- (log);
\end{tikzpicture}
\caption{Schema a blocchi del flusso sense-think-act-log in PomodorIA.}
\label{fig:diagramma-blocchi}
\end{figure}

\subsection{Mapping sui sette livelli IoTWF}

La Tabella~\ref{tab:iotwf} mappa esplicitamente ogni livello del modello
IoT World Forum (Sezione~\ref{sec:iotwf}) sul relativo componente software
del progetto, distinguendo fra ciò che accadrebbe in un deployment reale su
Raspberry Pi e ciò che è stato effettivamente implementato nel PoC.

\begin{table}[h]
\centering
\small
\begin{tabular}{@{}p{2.6cm}p{5.6cm}p{5.6cm}@{}}
\toprule
\textbf{Livello IoTWF} & \textbf{Deployment reale (target)} & \textbf{Implementazione nel PoC} \\
\midrule
1. Physical Devices \& Controllers &
Fotocamera, sensori DHT22/igrometro, pompa, ventola &
Moduli Python mock: \texttt{VirtualCameraSensor}, \texttt{EnvironmentSensorSimulator}, \texttt{FakeGPIO} \\
2. Connectivity &
Wi-Fi del Raspberry Pi &
Comunicazione in-memory fra moduli Python \\
3. Edge (Fog) Computing &
Inferenza CNN quantizzata sul Raspberry Pi &
Inferenza CNN compressa su CPU, con vincoli di risorse simulati (1 core, no GPU) \\
4. Data Accumulation &
Storage locale su SD card &
Log strutturati su file CSV \\
5. Data Abstraction &
Normalizzazione dei dati per il livello applicativo &
Strutture dati tipizzate (\texttt{dataclass}): \texttt{SensorReading}, \texttt{InferenceResult}, \texttt{AgentDecision} \\
6. Application &
Dashboard di monitoraggio &
Dashboard interattiva real-time (Streamlit + Plotly) \\
7. Collaboration \& Processes &
Notifiche all'agricoltore, integrazione gestionale &
Notifiche e allarmi generati dall'agente e mostrati nell'interfaccia \\
\bottomrule
\end{tabular}
\caption{Mapping dell'architettura di PomodorIA sui sette livelli IoTWF.}
\label{tab:iotwf}
\end{table}

\subsection{Le due modalità di esecuzione}

Il ciclo sense-think-act-log può essere eseguito da due programmi
alternativi, che orchestrano esattamente gli stessi cinque componenti di
dominio (fotocamera, sensori ambientali, motore di inferenza, agente
decisionale, banco attuatori), differenziandosi solo per l'interfaccia con
l'utente e per la modalità di persistenza dei risultati.

\begin{table}[h]
\centering
\small
\begin{tabular}{@{}p{3.4cm}p{5.6cm}p{5.6cm}@{}}
\toprule
& \textbf{\texttt{orchestrator/main\_loop.py}} & \textbf{\texttt{dashboard/app\_streamlit.py}} \\
\midrule
Interfaccia & Riga di comando (terminale) & Interfaccia web interattiva (browser) \\
Meccanismo del ciclo & Ciclo infinito nativo Python (\texttt{while True}) & Un ciclo per ogni \emph{rerun} di Streamlit, pilotato da un timer software \\
Persistenza & File CSV strutturato in \texttt{logs/} & Cronologia mantenuta in memoria (\texttt{st.session\_state}), persa alla chiusura \\
Caso d'uso tipico & Raccolta dati per analisi statistica offline (es. su molti cicli) & Dimostrazione interattiva e monitoraggio in tempo reale (es. in sede di esame) \\
\bottomrule
\end{tabular}
\caption{Confronto fra le due modalità di esecuzione del sistema.}
\label{tab:driver-confronto}
\end{table}

Questa duplicazione non introduce alcuna duplicazione di logica applicativa:
entrambi i programmi si limitano a istanziare gli stessi cinque componenti e
a richiamarli nello stesso ordine, garantendo che il comportamento del
sistema osservato da terminale e quello osservato dalla dashboard siano
concettualmente equivalenti.
